\paragraph{Decisiones que quedan para Lucía.}
\begin{itemize}
\item \emph{Resuelta:} \etq{P-03}, el camino \texttt{freeslip} plano. Corregido y medido
  (\etq{D-50}, §\ref{sec:defecto}).
\item \textbf{Orden y paso de tiempo si se usa la superficie deformable en producción}
  (§\ref{sec:celda}). No hace falta para $\alpha$ a orden lineal.
\end{itemize}

\paragraph{Lo que el modelo no hace.}
\begin{itemize}
\item \textbf{La geometría no lineal de la superficie.} El modelo es lineal en la frontera:
  la jerarquía cúbica de la sesión anterior
  (\archivo{research/2026-09-25_nonlinear_free_surface/}) es el paso siguiente natural,
  porque resuelve este mismo operador lineal una vez por orden, con los términos de orden
  superior como dato. Más allá de pocos órdenes hace falta una coordenada que siga la
  superficie (ALE), con operadores de coeficientes variables en todo el volumen: es un
  proyecto de otra escala. Un grafo tampoco representa olas que rompen.
\item \textbf{Validación a resolución de producción.} La puerta usa $8\times8\times128$; la
  estabilidad con capilaridad en $k$ grande, \texttt{ORD=4}, y más de dos procesos MPI no se
  probaron.
\item \textbf{Reinicio.} $\eta$ se escribe y se relee (\archivo{fs_eta.NNNN.dat}); E2 lo
  prueba (§\ref{sec:extras}). La primera corrida de un reinicio recalcula $B_{\mathbf k}$ y
  las trazas a partir del campo leído.
\end{itemize}

\paragraph{Problemas encontrados en el camino.}
\begin{itemize}
\item La referencia de Chebyshev de cuarto orden empeoraba al refinar y se reemplazó por la
  MAC de segundo orden con Richardson.
\item Un chequeo del límite de tapa rígida ($g\to\infty$ contra \texttt{freeslip}) se
  descartó antes de congelar la puerta: arrancando con $\eta=0$ se excitan ondas cuyo efecto
  escala como $g^{-1/2}$ y no como $g^{-1}$, así que el criterio habría sido frágil. Se
  reemplazó por F7, que es exacto.
\item El aborto de SPECTER upstream ante una cadena desconocida termina en un
  \emph{segmentation fault}, probablemente porque llama
  \texttt{MPI\_FINALIZE(MPI\_COMM\_WORLD, ierr)}, con un argumento de más (inferencia, no
  verificada). La puerta de la Fase~2 lo acepta porque sólo pide que aborte. No se tocó; el
  código nuevo aborta con \texttt{MPI\_ABORT}.
\item La hipótesis de la condición inicial incompatible como causa del F6 fallido era
  razonable y resultó falsa; se refutó con la referencia antes de tocar código.
\item La primera corrección del defecto (reconstruir sólo el incremento del subpaso) no
  cambió nada y se sacó del código.
\end{itemize}
